%%=============================================================================
%% CAPITULO 1 -- INTRODUCAO
%%
%% Na introducao, faca a apresentacao do trabalho. Ela devera ocupar uma ou,
%% no maximo, duas paginas. De um panorama da tematica, exponha a questao
%% problema que conduzira a pesquisa proposta, a hipotese e a justificativa.
%%
%% Titulos: nao utilize ponto apos a numeracao das secoes.
%%
%% Texto dummy (exercicio ia-6.1): trechos do pre-projeto do autor (versao 7).
%%=============================================================================

\section{Introdução}

A transformação digital da Administração Pública e a crescente dependência de
infraestruturas tecnológicas ampliaram a importância da gestão de riscos
cibernéticos para a continuidade de serviços, a proteção de informações e a
capacidade institucional de cumprir sua missão. Ao mesmo tempo, a expansão de
ambientes digitais aumenta a quantidade e a complexidade dos dados produzidos por
processos de segurança da informação, entre eles registros de ativos,
vulnerabilidades, exposições, criticidade e indicadores de ameaças. Esses dados
possuem elevado potencial para apoiar pesquisa, experimentação, auditoria,
desenvolvimento de modelos analíticos e avaliação de políticas e controles, mas
também podem revelar características sensíveis da infraestrutura e da capacidade
defensiva de uma organização.

Essa tensão produz um problema relevante para a Ciência de Dados aplicada ao setor
público: modelos analíticos dependem de dados suficientemente ricos para capturar
relações entre variáveis, enquanto a organização possui razões legítimas para
restringir a exposição dos dados originais. A geração de dados sintéticos surge,
nesse contexto, como uma alternativa para ampliar o acesso controlado a dados de
pesquisa e experimentação, sem pressupor que a simples substituição de registros
reais por registros artificiais seja suficiente para preservar o valor analítico.

Assim, estabelece-se como problema de pesquisa a necessidade de produzir dados
sintéticos de risco cibernético que sejam suficientemente fiéis e úteis para
análise, sem pressupor que a geração sintética, por si só, assegure privacidade ou
validade analítica.
